First suppose one reward location is known and the other is uniformly among $u$ unknown cells, with $t\ge1$ moves left and enough unknown cells. The optimal success probability is $\dfrac{2t-1}{u}$: inspect two new cells in each of the first $t-1$ moves until the other reward is found, then pair the known rewards next move; if none is found, on the last move pair the known reward with one new cell. At most $2t-1$ candidates can lead to a win, since inspecting two new cells on the last move cannot win when the known reward is elsewhere. Now let $F(u,t)$ count winning unordered placements among $\binom u2$ when neither reward is known. A first move should inspect two new cells: an already known empty cell yields no information or chance of winning. There is one placement with both rewards in those cells. For $t\ge2$, placements with one reward there contribute $2(2t-3)$ winning possibilities by the preceding bound, and placements with neither contribute $F(u-2,t-1)$. Thus $F(u,t)=1+2(2t-3)+F(u-2,t-1)$, with $F(u,1)=1$. It follows that $F(20,5)=1+3+7+11+15=37$, and the optimal probability is $\boxed{\dfrac{37}{190}}$.
